% =====================================================================
% TwoRoom appendix -- exact configuration, hyperparameters, protocol.
% Every value transcribed from tworoom_matrix_20260727/HYPERPARAMS.md,
% which was itself read out of the running configs / checkpoint
% config.json rather than from notes.
%
% Drop in after the main body. If the paper already opens an appendix
% (for app:proofs), delete the \appendix line below.
% Requires \usepackage{booktabs} (already in the preamble).
% =====================================================================

\appendix

\section{TwoRoom: exact configuration and evaluation protocol}
\label{app:tworoom}

This appendix specifies everything needed to reproduce
Section~\ref{sec:tworoom}. Values were read out of the running
configuration and the checkpoint \texttt{config.json} files.

\subsection{Environment}

\begin{table}[h]
\centering\small
\begin{tabular}{ll}
\toprule
observation & $224\times224$ RGB, pure CPU rasteriser (no physics engine) \\
action space & $[-1,1]^2$, interpreted as a velocity direction \\
max speed & $10.5$ px/step \\
control rate & $10$ Hz; frameskip $5$ (one action block $=5$ primitive steps) \\
geometry & image $224$, border $14$, vertical wall at $x = 112$ \\
doorways & $1$--$3$, count and positions re-randomised every episode \\
success & agent centre within $16$ px of the target; episode terminates \\
randomised per reset & agent position, target position (target in the opposite room) \\
episode cap & set by the plan config to $2\times$ the evaluation budget \\
\bottomrule
\end{tabular}
\end{table}

\subsection{Offline dataset}

A single demonstration set is used for pretraining the world models, for
fitting every learned component, and for constructing the evaluation
tasks (see~\ref{app:protocol}).

\begin{table}[h]
\centering\small
\begin{tabular}{ll}
\toprule
frames / episodes & $920{,}809$ / $10{,}000$ ($\approx 92$ frames per episode) \\
collection policy & the environment's scripted weak expert (no action noise, no repeat) \\
columns used & \texttt{pixels} (uint8), \texttt{action} $\in\mathbb R^2$,
                \texttt{proprio} $\in\mathbb R^2$, \texttt{pos\_agent} $\in\mathbb R^2$ \\
action representation & $10$-d $z$-scored \emph{blocks} $=$ $2$-d action $\times$ frameskip $5$ \\
action statistics & $z$-scoring fit on this same file, matching the pretraining convention \\
\bottomrule
\end{tabular}
\end{table}

\subsection{World models (frozen throughout)}

All three are author-released checkpoints, used frozen; only the critic
and the planner are trained here. LeWM and PLDM are architecturally
identical---verified field-by-field from both configs---and differ only in
pretraining objective, which makes them a controlled pair.

\begin{table}[h]
\centering\small
\begin{tabular}{lll}
\toprule
 & LeWM \,/\, PLDM & DINO-WM \\
\midrule
encoder & ViT-tiny, patch $14$, image $224$, trained from scratch & frozen \texttt{dinov2\_small} \\
latent & single $192$-d \texttt{CLS} token & $196$ patch tokens, $77{,}224$-d flat \\
token composition & --- & [pixel $384 \,\Vert\,$ proprio-emb $10$] $= 394$ \\
predictor & adaLN, depth $6$, heads $16$, mlp $2048$, head-dim $64$ & causal, same depth/heads/mlp \\
predictor width & $192$ & $404$ ($=394 + $ action-emb $10$) \\
action encoder & linear $10 \to 192$ & linear $10 \to 10$ \\
proprio & \textbf{not consumed} & linear $2 \to 10$, image resized $224\to196$ \\
history & $3$ frames & $3$ frames ($588$ tokens per forward) \\
objective & $\mathcal L_{\text{pred}} + \lambda\mathcal L_{\text{SIGReg}}$ \;/\;
            sim $+$ std $+$ cov $+$ temporal $+$ IDM & --- \\
\bottomrule
\end{tabular}
\end{table}

Conversion of the released checkpoints into our runner was validated
before any planning: encoder outputs match the reference implementation to
$2.4$--$2.8\times10^{-5}$ max absolute difference, and one-step predictor
MSE relative to a copy-last baseline is $0.043$ (LeWM), $0.168$ (PLDM),
$0.102$ (DINO-WM), with shuffled-action ratios strictly above the
copy-last ratios in all three cases, confirming genuine action
conditioning rather than a degenerate predictor.

\subsection{Evaluation protocol}
\label{app:protocol}

Tasks are constructed by replaying offline trajectories, which fixes the
task distribution to the one the world models were pretrained on and makes
the goal distance an explicit control variable:

\begin{enumerate}\itemsep2pt
  \item Draw a row $t$ from the dataset (the seed drives the draw).
  \item Reset the environment to the recorded agent position at row $t$.
  \item Set the goal to the recorded agent position at row $t + \Delta$;
        the goal \emph{image} handed to the planner is that row's frame.
  \item Roll out the receding-horizon controller until success ($16$ px)
        or until the budget expires.
\end{enumerate}

\begin{table}[h]
\centering\small
\begin{tabular}{ll}
\toprule
goal distance & $\Delta = 25$ primitive steps ($h25$, the published
                setting) and $\Delta = 50$ ($h50$, our extension) \\
budget & $2\Delta$ primitive steps ($50$ / $100$): twice the length of the
         demonstrated solution \\
planning cycles & $2$ at $h25$, $4$ at $h50$ ($25$ primitive steps executed
                  per replan) \\
episodes per cell & $50$ \\
task seeds & $42, 43, 44$ (drive the task draw) \\
initialisation seeds & $0, 1, 2$ for every learned component
                       (DINO-WM planner: $0$ only) \\
metric & success rate; a reported cell is the mean over $3$ (untrained) or
         $9$ (learned) runs \\
controller & horizon $N = 5$ action blocks, replan interval $5$ blocks
             $=25$ primitive steps, $10$ environments per solver call \\
\bottomrule
\end{tabular}
\end{table}

Two references calibrate the numbers. Published $h25$ success for these
checkpoints is $87$ (LeWM), $97$ (PLDM), $100$ (DINO-WM), which our
protocol reproduces at $89.3$, $96.7$, $100.0$. Under the identical
protocol, a policy that never moves scores $8.0$ and a uniform-random
policy $28.0$, which is the floor the reactive-policy row of
Table~\ref{tab:tworoom} should be read against.

\subsection{Planners}

\begin{table}[h]
\centering\small
\begin{tabular}{lll}
\toprule
planner & setting & value \\
\midrule
CEM  & candidates / iterations / elites / var & $300$ / $30$ ($10$ at the
       published budget) / $30$ / $1.0$ \\
MPPI & candidates / iterations / elites / temperature & $300$ / $30$ /
       $30$ / $0.5$ \\
Adam & candidates / gradient steps / optimiser & $100$ / $30$ /
       AdamW, lr $0.1$, no action noise \\
\midrule
RLP  & refinement steps $K$ & $8$ (reported), truncated to $4$ and $2$ at deployment \\
     & action clip $a_{\max}$ / max residual & $2.2$ / $12$ \\
     & width / layers / heads / feed-forward / $z$-projection & $256$ / $2$ / $4$ / $2\times$width / $64$ \\
     & goal parameterisation & displacement, $\mathrm{proj}(\hat z_g - \hat z_t)$ \\
     & deployment & no restarts, no sampling, no robust averaging \\
\bottomrule
\end{tabular}
\end{table}

Three caveats belong with this table. \textbf{(i)} No plan-time MPPI
configuration ships with the reference implementation; ours mirrors the CEM
config and takes the class-default temperature, and temperature---the
parameter MPPI is most sensitive to---was never swept, so those rows are a
lower bound and are not used to support any claim. \textbf{(ii)} Adam is
the only planner that retains a backward graph over all candidates. On
DINO-WM's $588$-token predictor this costs ${\approx}0.27$ GB per retained
candidate on top of a ${\approx}25$ GB base, so $200$ candidates peak at
$78.5$ GB and do not fit an $80$ GB card; DINO-WM's Adam arm therefore ran
one environment per solver call instead of ten. That is environment
chunking and does not change the algorithm, but it does change the random
realisation, so those two cells are statistically equivalent to---not a
bit-identical protocol with---the others. \textbf{(iii)} The implemented
refiner conditions on slightly more than the interface stated in
Section~\ref{sec:tworoom}: each plan token carries
$[\mathbf a_k^{(t)},\, \nabla_{\!\mathbf a_k^{(t)}} V,\,
\mathrm{proj}(\hat z_t),\, \mathrm{proj}(\hat z_g - \hat z_t)]$ plus a
positional and a per-iteration embedding, i.e.\ it also sees the current
and goal latents, not only the scalar value and its gradient.

\subsection{Critic (the value metric)}

The critic is a metric residual network, i.e.\ a learned quasimetric
\begin{equation*}
  d_\psi(\hat z_i \to \hat z_j)
  = \lVert u(\hat z_i) - u(\hat z_j) \rVert_2
  \;+\; \max_k \mathrm{ReLU}\bigl( v(\hat z_j)_k - v(\hat z_i)_k \bigr),
\end{equation*}
where $u$ and $v$ are the symmetric and asymmetric halves of a single
$128$-d embedding produced by an MLP $d_z \to 256 \to 256 \to 128$ with
SiLU activations. It is non-negative and obeys the triangle inequality by
construction, so long distances compose out of short transitions.

\begin{table}[h]
\centering\small
\begin{tabular}{ll}
\toprule
symmetric fraction & $0.5$ ($64$ symmetric $\Vert$ $64$ asymmetric); fixed, never varied \\
target & $n$-step, $n = 50$ primitive steps: Monte-Carlo if the goal is
         reached within $n$, else bootstrap \\
discount & $1.0$ (the target is true steps-to-go) \\
loss & expectile Huber, expectile $0.1$, $\beta = 1.0$ \\
hindsight goals & $0.3$ cross-episode (for stitching), otherwise
                  in-episode at a bucket-balanced offset ($10$ buckets) \\
steps / batch & $6{,}000$ / $1{,}024$ \\
optimiser & AdamW, lr $10^{-3}$, weight decay $10^{-4}$; target network
            Polyak $\tau = 0.005$ \\
seeds / wall-clock & $0,1,2$; ${\approx}90$ s per seed ($192$-d),
                     ${\approx}22$ min ($77{,}224$-d) \\
\bottomrule
\end{tabular}
\end{table}

A low expectile is deliberate: for a cost-to-go, optimism should point
toward the minimum over actions. One approximation is worth stating
explicitly---the critic is fit on triples drawn entirely from encoder
latents, but at plan time it scores \emph{predictor} outputs against an
encoder goal. It is trained encoder-to-encoder and deployed
predictor-to-encoder, which is benign only insofar as the predictor's
output stays in encoder space.

\subsection{Planner training}

\begin{table}[h]
\centering\small
\begin{tabular}{ll}
\toprule
critic initialisation & warm-started from the TD checkpoint \emph{of the same seed} \\
critic during planner training & co-trained, expectile annealed
                                $0.1 \to 0.03$, lr $10^{-3} \to 10^{-4}$ \\
planner lr & $3\times10^{-4} \to 3\times10^{-5}$ \\
steps / batch & $8{,}000$ / $128$ (LeWM, PLDM); $8$ at $K{=}8$ and $16$ at
                $K{=}4$ for DINO-WM (memory) \\
$n$-step / $K$ / $a_{\max}$ / max residual & $50$ / $8$ / $2.2$ / $12$ \\
seeds / wall-clock & $0,1,2$ (DINO-WM: $0$); ${\approx}58$ min per seed
                     ($192$-d), ${\approx}4$ h ($77{,}224$-d) \\
\bottomrule
\end{tabular}
\end{table}

Note the deviation from Section~\ref{sec:theoretical-results}, which for
clarity treats the critic as fixed while the planner trains: in the runs
reported here the critic is warm-started from the TD solution and then
co-trained with the planner, with the expectile annealed downward. The
optimality statement is unaffected---it is conditional on the critic---but
the empirical numbers correspond to the co-trained variant. A separate
action-clip sweep ($a_{\max} \in \{2.2, 2.8, 3.5\}$, $K \in \{8,12\}$)
selected $a_{\max} = 2.8$ for PLDM on task seed $42$ alone; every number in
Table~\ref{tab:tworoom} uses the unswept canonical $a_{\max} = 2.2$, so no
reported cell benefits from that selection.

\subsection{Reactive-policy baseline}

The single-pass baseline is
$\pi(\hat z, \hat z_g) = \tanh\bigl(\mathrm{MLP}[\mathrm{proj}(\hat z)
\,\Vert\, \mathrm{proj}(\hat z_g - \hat z)]\bigr)\cdot a_{\max}$ --- the
planner's goal parameterisation and action box, with the refinement
recursion removed. It is trained by first-order extraction through the
frozen world model, in the manner of PWM \citep{georgiev2024pwm} but fully
offline, and scored by the same quasimetric critic rather than a learned
reward model and critic ensemble.

\begin{table}[h]
\centering\small
\begin{tabular}{ll}
\toprule
objective & $J = -\sum_{t=1}^{H} \gamma^t\, d_{\bar\psi}(\hat z_t, \hat z_g)$,
            loss $\mathbb E[-J/H]$, latents imagined through the frozen model \\
horizon / discount / action clip & $5$ / $0.99$ / $2.2$ \\
actor & width $512$, $3$ layers, $z$-projection $256$; AdamW
        $5\times10^{-4} \to 5\times10^{-5}$, gradient norm $100$ \\
critic / teacher & same MRN, warm-started and co-trained; Polyak EMA
                   teacher, $\tau = 0.005$ \\
sampler & $n = 50$, cross-episode goal probability $0.3$, offsets capped at
          $12$ blocks, bucket-balanced \\
steps / batch / seeds & $8{,}000$ / $128$ (DINO-WM: $16$) / $0,1,2$
                        (DINO-WM: $0$) \\
deployment & one forward pass per replan; the plan tail is either imagined
             or held, and the two agree to $\le 5.6$ points \\
\bottomrule
\end{tabular}
\end{table}

The discount, actor learning rate and gradient clipping follow PWM; what is
replaced is the reward model and three-critic ensemble (by the quasimetric
and its EMA teacher) and the online interaction (by the offline cache).
The world model stays frozen, whereas PWM fine-tunes it. This baseline is
therefore evidence about single-pass extraction \emph{on this stack}, not a
reproduction of PWM in its own setting; it received no hyperparameter
search.

\subsection{Latent caches and compute}

Learned components are fit on cached latents rather than raw pixels. Two
strides are used: every frame for critic training, so that $n$-step targets
see every transition, and every fifth frame---one latent per action
block---for planner rollout contexts. The single-token caches are $696$ MB
each. DINO-WM at full width over all frames would be $284$ GB, so its cache
is capped at $200{,}000$ rows ($58$ GB) at full latent width. The cap takes
a contiguous prefix, i.e.\ roughly episodes $0$--$2{,}170$, so the DINO-WM
learned rows are fit on about a fifth of the corpus. A $1024$-d random
projection is $80\times$ smaller and works for the critic, whose wrapper
dispatches on width, but not for the planner, which bakes the latent width
into its input projections and would then shape-mismatch the full-width
latents supplied at deployment.

All experiments ran on $4\times$H100 80 GB. A CEM evaluation cell costs
${\approx}34$ min on DINO-WM and a few minutes on the single-token models;
the Adam cells on DINO-WM cost ${\approx}2.5$ h each.

\subsection{Reproduction notes}

Three implementation details changed measured outcomes materially and are
recorded so that a reimplementation can check for them.

First, \textbf{double normalisation of proprioception}. Our runner
$z$-scores every cached observation key, and DINO-WM's encoder then applies
its own proprioception statistics, fit from the same data. Applying both
maps every agent position to ${\approx}-3\sigma$: positions $2$ px apart
arrive identical to three decimals, destroying the goal's position signal.
Passing proprioception through unnormalised on that model alone moves it
from $36.0$ to $100.0$ at $h25$ --- from far below its published value to
matching it exactly --- and also makes the run ${\approx}1.6\times$ faster,
because the planner converges instead of thrashing. Only DINO-WM consumes
proprioception, so the other two are unaffected.

Second, \textbf{CPU thread caps are load-bearing}. The renderer is
CPU-bound; leaving BLAS thread counts uncapped slowed evaluation by roughly
$40\times$. We cap them at $12$--$16$.

Third, \textbf{the encoder checkpoint layout is library-version
dependent}. The released single-token encoders use a pre-$5$.x Hugging Face
ViT parameter layout; converting them under the wrong major version has
\texttt{load\_state\_dict} reject all $192$ encoder tensors. Our converter
detects the installed version. Reference environment: torch
\texttt{2.4.1+cu124}, transformers \texttt{4.57.6}, Python \texttt{3.11.10}.
